\section{Related work}
\label{sec:related}

\paragraph{A note on how these references are used.}
Every reference below was checked against the publisher or author page for
authors, title, venue and year, and the one-line method descriptions are at the
level the titles and abstracts support. No number in this paper is attributed to
any of them. We compare against our own analytic baseline and against ExpandNet
run from its authors' released weights (\cref{sec:expandnet});
\cref{sec:limitations} states which comparisons are still missing. The failure
mode this caution replaces, placeholder tables presented as measurements, is
documented in \cref{sec:withdrawal}.

\paragraph{Single-image inverse tone mapping.}
The dominant framing learns a mapping from an 8-bit frame to a higher-range one,
typically with an encoder-decoder and a loss in a perceptual or log domain, and
typically evaluated on synthetically clipped data. Eilertsen et
al.~\citep{eilertsen2017} predict a log-domain reconstruction of the saturated
regions with a U-Net and blend it back through a saturation mask; Marnerides et
al.~\citep{marnerides2018} combine three branches at different receptive fields;
Endo et al.~\citep{endo2017} instead synthesise a bracketed exposure stack from
the single input and merge it; Liu et al.~\citep{liu2020} decompose the task into
learned inverses of the individual camera-pipeline stages. Our architecture is
deliberately smaller than this line of work and is a residual on an analytic
inverse tone map rather than a direct predictor, which is what makes ``does the
network beat doing nothing?'' a question we can ask at every step of training.

\paragraph{Highlight and clipped-region reconstruction.}
A related family treats the problem as inpainting the saturated regions
specifically, rather than remapping the whole frame. Santos et
al.~\citep{santos2020} make the clipped area explicit by masking it out of the
network's own features and adding a perceptual loss, and the masked blend of
Eilertsen et al.~\citep{eilertsen2017} serves the same end. Our highlight and
shadow masks share that purpose but are predicted jointly with the residual and
gated by a luminance prior rather than by a detected saturation mask.
\Cref{sec:failure} is a direct criticism of that choice: our own failures are
not in the clipped highlights at all.

\paragraph{Evaluation of HDR reconstruction.}
PU21 \citep{pu21} and ColorVideoVDP \citep{cvvdp} are the two public measuring
sticks we use, the first an encoding that makes existing metrics such as PSNR
applicable to absolute-luminance HDR, the second a calibrated visible-difference
predictor reporting in JOD. \Cref{sec:results} reports a case where they diverge sharply in the practical
significance they assign to the same difference on the same frames, and
\cref{sec:failure} shows the PU21 penalty is only partly tail-carried. We are
not aware of that disagreement having been characterised for this task, and
consider it a contribution independent of the model.

\paragraph{Training data for HDR.}
The scarcity of paired SDR/HDR footage shapes every result here. Public HDR
video with cinematic production values remains small enough to enumerate: the
HDM-HDR-2014 set of Froehlich et al.~\citep{froehlich2014} contributes 9 scenes
and \num{11007} frames to our corpus, and the widely used HDR still sets
\citep{kalantari2017} are bracketed-exposure captures rather than graded footage.
The remainder of our corpus is proprietary. \Cref{sec:corpus} documents a
distribution shift between our own training and held-out splits that we did not
design and that a reader should weigh.

\paragraph{What we compare against.}
ExpandNet \citep{marnerides2018} is scored on our split in \cref{sec:expandnet},
using the authors' own model and released weights and their own pre- and
post-processing. Our benchmark harness scores any third-party output against the
same reference, so the remaining comparisons are a matter of running published
code rather than of building anything; \cref{sec:limitations} says which are
outstanding.
